% GENERATED FILE. Edit canonical lessons, then run npm run generate:books.
% canonical-source-hash: fnv1a64:cdbb7674fd66e7d3
% canonical-lessons: UR-C291-hilna, UR-C291-jhatka, UR-C291-darar, UR-C291-dhansna, UR-C291-dabna, UR-R291-review-284-291

\chapter{When the Ground Shook}
\label{ch:ur-when-the-ground-shook}
\hlchaptermodality{291}

\begin{chapteropening}
\textbf{By the end of this chapter:} \emph{I can tell what happened in an earthquake.}

Complete the last lesson of chapter 291: I can tell what happened in an earthquake.
\end{chapteropening}

\section[\texorpdfstring{\ur{ہلنا} (hilnā)}{ہلنا (hilnā)}]{\ur{ہلنا} (hilnā) — to shake (by itself); to move}

\label{lesson:UR-C291-hilna}



\begin{quote}
Before the new one: say the Urdu for to be averted, then the Urdu for caution.
\end{quote}

\subsection*{You'll want to know: \ur{ہلنا}}
\textbf{\ur{ہلنا}} — \emph{hilnā} — \textquotedblleft{}to shake (by itself); to move\textquotedblright{}.

\begin{quote}
\ur{رات} \ur{کو} \ur{زمین} \ur{ہلی۔}
\end{quote}

\begin{quote}
\emph{rāt ko zamīn hilī.}
\end{quote}

\begin{quote}
The ground shook in the night.
\end{quote}

Past (it …): \ur{وہ} \ur{زمین} \ur{ہلی}, \emph{voh zamīn hilī}.

\begin{cousinweb}
\textbf{\ur{ہلنا}} (\emph{hilnā}) is inherited: CDIAL lists the Hindi cousin \emph{hilnā}, \textquotedblleft{}to shake\textquotedblright{}, under a form it reconstructs (CDIAL 14120). It is a thing moving by itself: the ground, a tooth, a leaf.
\end{cousinweb}

\begin{grammarlens}[title={What you've built}]
One more word for telling what happened.
\end{grammarlens}

\subsection*{Your turn}
\begin{itemize}
\raggedright
  \item \emph{Say it:} \emph{hilnā}
  \item \emph{Say it:} \emph{hilnā}, once more
  \item \emph{Say it:} \emph{rāt ko zamīn hilī.} — The ground shook in the night.
  \item \emph{Recall:} say the Urdu for to be averted, then the Urdu for caution, then say \emph{hilnā} again
\end{itemize}

\begin{tcolorbox}[breakable,colback=teal!4,colframe=teal!35!black,title={Before you move on}]
What does \ur{ہلنا} mean? (\textquotedblleft{}To shake (by itself); to move\textquotedblright{}.) Say it once more.
\end{tcolorbox}
\section[\texorpdfstring{\ur{جھٹکا} (jhaṭkā)}{جھٹکا (jhaṭkā)}]{\ur{جھٹکا} (jhaṭkā) — a jolt; a tremor}

\label{lesson:UR-C291-jhatka}



\begin{quote}
Before the new one: say the Urdu for caution, then the Urdu for to shake (by itself).
\end{quote}

\subsection*{You'll want to know: \ur{جھٹکا}}
\textbf{\ur{جھٹکا}} — \emph{jhaṭkā} — \textquotedblleft{}a jolt; a tremor\textquotedblright{}.

\begin{quote}
\ur{پہلا} \ur{جھٹکا} \ur{ہلکا} \ur{تھا۔}
\end{quote}

\begin{quote}
\emph{pehlā jhaṭkā halkā thā.}
\end{quote}

\begin{quote}
The first jolt was mild.
\end{quote}

Masculine.

\begin{cousinweb}
\textbf{\ur{جھٹکا}} (\emph{jhaṭkā}) is inherited: CDIAL lists the Hindi cousin \emph{jhaṭkā}, \textquotedblleft{}a jerk\textquotedblright{}, under \emph{*jhaṭ}, a form it reconstructs (CDIAL 5327).
\end{cousinweb}

\begin{grammarlens}[title={What you've built}]
One more word for telling what happened.
\end{grammarlens}

\subsection*{Your turn}
\begin{itemize}
\raggedright
  \item \emph{Say it:} \emph{jhaṭkā}
  \item \emph{Say it:} \emph{jhaṭkā}, once more
  \item \emph{Say it:} \emph{pehlā jhaṭkā halkā thā.} — The first jolt was mild.
  \item \emph{Recall:} say the Urdu for caution, then the Urdu for to shake (by itself), then say \emph{jhaṭkā} again
\end{itemize}

\begin{tcolorbox}[breakable,colback=teal!4,colframe=teal!35!black,title={Before you move on}]
What does \ur{جھٹکا} mean? (\textquotedblleft{}A jolt; a tremor\textquotedblright{}.) Say it once more.
\end{tcolorbox}
\section[\texorpdfstring{\ur{دراڑ} (darāṛ)}{دراڑ (darāṛ)}]{\ur{دراڑ} (darāṛ) — a crack}

\label{lesson:UR-C291-darar}



\begin{quote}
Before the new one: say the Urdu for to shake (by itself), then the Urdu for a jolt.
\end{quote}

\subsection*{You'll want to know: \ur{دراڑ}}
\textbf{\ur{دراڑ}} — \emph{darāṛ} — \textquotedblleft{}a crack\textquotedblright{}.

\begin{quote}
\ur{دیوار} \ur{میں} \ur{دراڑ} \ur{آ} \ur{گئی۔}
\end{quote}

\begin{quote}
\emph{dīvār meṅ darāṛ ā gaī.}
\end{quote}

\begin{quote}
A crack appeared in the wall.
\end{quote}

Feminine.

\begin{cousinweb}
\textbf{\ur{دراڑ}} (\emph{darāṛ}) is inherited: CDIAL lists the Hindi cousin \emph{darār}, \textquotedblleft{}a crack, a rent\textquotedblright{}, under a form it reconstructs (CDIAL 6193); Urdu writes it with \textbf{\ur{ڑ}}. A crack is said to come: \textbf{\ur{دراڑ} \ur{آ} \ur{گئی}}.
\end{cousinweb}

\begin{grammarlens}[title={What you've built}]
One more word for telling what happened.
\end{grammarlens}

\subsection*{Your turn}
\begin{itemize}
\raggedright
  \item \emph{Say it:} \emph{darāṛ}
  \item \emph{Say it:} \emph{darāṛ}, once more
  \item \emph{Say it:} \emph{dīvār meṅ darāṛ ā gaī.} — A crack appeared in the wall.
  \item \emph{Recall:} say the Urdu for to shake (by itself), then the Urdu for a jolt, then say \emph{darāṛ} again
\end{itemize}

\begin{tcolorbox}[breakable,colback=teal!4,colframe=teal!35!black,title={Before you move on}]
What does \ur{دراڑ} mean? (\textquotedblleft{}A crack\textquotedblright{}.) Say it once more.
\end{tcolorbox}
\section[\texorpdfstring{\ur{دھنسنا} (dhansnā)}{دھنسنا (dhansnā)}]{\ur{دھنسنا} (dhansnā) — to cave in; to sink in}

\label{lesson:UR-C291-dhansna}



\begin{quote}
Before the new one: say the Urdu for a jolt, then the Urdu for a crack.
\end{quote}

\subsection*{You'll want to know: \ur{دھنسنا}}
\textbf{\ur{دھنسنا}} — \emph{dhansnā} — \textquotedblleft{}to cave in; to sink in\textquotedblright{}.

\begin{quote}
\ur{سڑک} \ur{بیچ} \ur{میں} \ur{دھنس} \ur{گئی۔}
\end{quote}

\begin{quote}
\emph{saṛak bīch meṅ dhans gaī.}
\end{quote}

\begin{quote}
The road caved in in the middle.
\end{quote}

Past (it …): \ur{وہ} \ur{سڑک} \ur{دھنس} \ur{گئی}, \emph{voh saṛak dhans gaī}.

\begin{cousinweb}
\textbf{\ur{دھنسنا}} (\emph{dhansnā}) is inherited: CDIAL lists the Hindi cousin \emph{dhãsnā}, \textquotedblleft{}to sink into\textquotedblright{}, under Sanskrit \emph{dhváṃsati} (CDIAL 6896).
\end{cousinweb}

\begin{grammarlens}[title={What you've built}]
One more word for telling what happened.
\end{grammarlens}

\subsection*{Your turn}
\begin{itemize}
\raggedright
  \item \emph{Say it:} \emph{dhansnā}
  \item \emph{Say it:} \emph{dhansnā}, once more
  \item \emph{Say it:} \emph{saṛak bīch meṅ dhans gaī.} — The road caved in in the middle.
  \item \emph{Recall:} say the Urdu for a jolt, then the Urdu for a crack, then say \emph{dhansnā} again
\end{itemize}

\begin{tcolorbox}[breakable,colback=teal!4,colframe=teal!35!black,title={Before you move on}]
What does \ur{دھنسنا} mean? (\textquotedblleft{}To cave in; to sink in\textquotedblright{}.) Say it once more.
\end{tcolorbox}
\section[\texorpdfstring{\ur{دبنا} (dabnā)}{دبنا (dabnā)}]{\ur{دبنا} (dabnā) — to be pressed down; to be trapped under}

\label{lesson:UR-C291-dabna}



\begin{quote}
Before the new one: say the Urdu for a crack, then the Urdu for to cave in.
\end{quote}

\subsection*{You'll want to know: \ur{دبنا}}
\textbf{\ur{دبنا}} — \emph{dabnā} — \textquotedblleft{}to be pressed down; to be trapped under\textquotedblright{}.

\begin{quote}
\ur{کچھ} \ur{لوگ} \ur{چھت} \ur{کے} \ur{نیچے} \ur{دب} \ur{گئے۔}
\end{quote}

\begin{quote}
\emph{kuchh log chhat ke nīche dab gae.}
\end{quote}

\begin{quote}
Some people were trapped under the roof.
\end{quote}

Past (he …): \ur{وہ} \ur{دب} \ur{گیا}, \emph{voh dab gayā}.

\begin{cousinweb}
\textbf{\ur{دبنا}} (\emph{dabnā}) is inherited: CDIAL lists the Hindi cousin \emph{dabnā}, \textquotedblleft{}to be covered\textquotedblright{}, under \emph{*dabb}, a form it reconstructs (CDIAL 6173). It is being pressed under a weight.
\end{cousinweb}

\begin{grammarlens}[title={What you've built}]
One more word for telling what happened.
\end{grammarlens}

\subsection*{Your turn}
\begin{itemize}
\raggedright
  \item \emph{Say it:} \emph{dabnā}
  \item \emph{Say it:} \emph{dabnā}, once more
  \item \emph{Say it:} \emph{kuchh log chhat ke nīche dab gae.} — Some people were trapped under the roof.
  \item \emph{Recall:} say the Urdu for a crack, then the Urdu for to cave in, then say \emph{dabnā} again
\end{itemize}

\begin{tcolorbox}[breakable,colback=teal!4,colframe=teal!35!black,title={Before you move on}]
What does \ur{دبنا} mean? (\textquotedblleft{}To be pressed down; to be trapped under\textquotedblright{}.) Say it once more.
\end{tcolorbox}
\section[(dialogue)]{Looking back: from A Crash on the Road to When the Ground Shook}

\label{lesson:UR-R291-review-284-291}



\begin{quote}
Say the Urdu for speed and to collide.
\end{quote}

\subsection*{Your turn}
Say these aloud:

\begin{itemize}
\raggedright
  \item speed, to collide, a wheel, to overturn, a ride
  \item a stumble, to stagger, a foot, to be grazed, a scratch
  \item unconscious, covered in blood, medical treatment, to crush, to steady oneself
  \item an ember, to smoulder, a spark, to flare up, to be on fire
  \item an oil lamp, a flame, to glow red-hot, brightness, scorching heat
  \item a siren, a crew, to go out (of a fire), to melt, to rush in a panic
  \item a disaster, an announcement, a danger, to be averted, caution
  \item to shake (by itself), a jolt, a crack, to cave in, to be trapped under
\end{itemize}

\begin{tcolorbox}[breakable,colback=teal!4,colframe=teal!35!black,title={Before you move on}]
Speed? (\textbf{\ur{رفتار}}, \emph{raftār}.) To collide? (\textbf{\ur{ٹکرانا}}, \emph{ṭakrānā}.) A wheel? (\textbf{\ur{پہیا}}, \emph{pahiyā}.) To overturn? (\textbf{\ur{الٹنا}}, \emph{ulaṭnā}.) A ride? (\textbf{\ur{سواری}}, \emph{savārī}.) A stumble? (\textbf{\ur{ٹھوکر}}, \emph{ṭhokar}.) To stagger? (\textbf{\ur{لڑکھڑانا}}, \emph{laṛkhaṛānā}.) A foot? (\textbf{\ur{پاؤں}}, \emph{pāoṅ}.) To be grazed? (\textbf{\ur{چھلنا}}, \emph{chhilnā}.) A scratch? (\textbf{\ur{خراش}}, \emph{kharāsh}.) Unconscious? (\textbf{\ur{بے} \ur{ہوش}}, \emph{be hosh}.) Covered in blood? (\textbf{\ur{لہولہان}}, \emph{lahūluhān}.) Medical treatment? (\textbf{\ur{علاج}}, \emph{ilāj}.) To crush? (\textbf{\ur{کچلنا}}, \emph{kuchalnā}.) To steady oneself? (\textbf{\ur{سنبھلنا}}, \emph{sambhalnā}.) An ember? (\textbf{\ur{انگارا}}, \emph{angārā}.) To smoulder? (\textbf{\ur{سلگنا}}, \emph{sulagnā}.) A spark? (\textbf{\ur{شرارہ}}, \emph{sharārā}.) To flare up? (\textbf{\ur{بھڑکنا}}, \emph{bhaṛaknā}.) To be on fire? (\textbf{\ur{جلنا}}, \emph{jalnā}.) An oil lamp? (\textbf{\ur{چراغ}}, \emph{chirāgh}.) A flame? (\textbf{\ur{شعلہ}}, \emph{sholā}.) To glow red-hot? (\textbf{\ur{دہکنا}}, \emph{dahaknā}.) Brightness? (\textbf{\ur{روشنی}}, \emph{raushnī}.) Scorching heat? (\textbf{\ur{تپش}}, \emph{tapish}.) A siren? (\textbf{\ur{سائرن}}, \emph{sāiran}.) A crew? (\textbf{\ur{عملہ}}, \emph{amlā}.) To go out (of a fire)? (\textbf{\ur{بجھنا}}, \emph{bujhnā}.) To melt? (\textbf{\ur{پگھلنا}}, \emph{pighalnā}.) To rush in a panic? (\textbf{\ur{ہڑبڑانا}}, \emph{haṛbaṛānā}.) A disaster? (\textbf{\ur{آفت}}, \emph{āfat}.) An announcement? (\textbf{\ur{اعلان}}, \emph{elān}.) A danger? (\textbf{\ur{خطرہ}}, \emph{khatrā}.) To be averted? (\textbf{\ur{ٹلنا}}, \emph{ṭalnā}.) Caution? (\textbf{\ur{احتیاط}}, \emph{ehtiyāt}.) To shake (by itself)? (\textbf{\ur{ہلنا}}, \emph{hilnā}.) A jolt? (\textbf{\ur{جھٹکا}}, \emph{jhaṭkā}.) A crack? (\textbf{\ur{دراڑ}}, \emph{darāṛ}.) To cave in? (\textbf{\ur{دھنسنا}}, \emph{dhansnā}.) To be trapped under? (\textbf{\ur{دبنا}}, \emph{dabnā}.)
\end{tcolorbox}
